% concatenated from ferrers.tex
\documentclass[twoside, reqno]{amsart}

\usepackage[T1]{fontenc}
\usepackage[utf8]{inputenc}
\usepackage{amsmath,amssymb,amsthm}
\usepackage{microtype}
\usepackage[hyphens]{url}
\usepackage{hyperref}
\usepackage{enumitem}
\usepackage{xcolor}
\usepackage{graphicx}
\usepackage{caption}
\usepackage{mathtools}
\usepackage{thmtools}

\hypersetup{
  colorlinks=true,
  linkcolor=blue!50!black,
  citecolor=blue!50!black,
  urlcolor=blue!50!black,
}

\def\adj{\mathrel{\!\mathrel-\mkern-8mu\mathrel-\mkern-8mu\mathrel-\!}}
   % adjacent vertices, from taocpmac.tex

\newtheorem{theorem}{Theorem}
\newtheorem*{corollary}{Corollary}
\newtheorem{proposition}[theorem]{Proposition}
\newtheorem{lemma}[theorem]{Lemma}

\theoremstyle{definition}
\newtheorem{definition}[theorem]{Definition}
\newtheorem{example}[theorem]{Example}

\theoremstyle{remark}
\newtheorem{remark}[theorem]{Remark}

\newcommand{\RR}{\mathbb R}
\newcommand{\diag}{\operatorname{diag}}
\newcommand{\epsov}{\varepsilon}

% https://tex.stackexchange.com/questions/319398/making-the-dot-in-the-theorem-normal-but-not-as-like-bf
\makeatletter
\def\thmhead@plain#1#2#3{%
  \thmname{#1}\thmnumber{\@ifnotempty{#1}{ }\@upn{#2}}%
  \thmnote{ \the\thm@notefont(#3)}}% no group inside \thmnote
\let\thmhead\thmhead@plain
\makeatother

\title[The Ferrers bound for spanning trees in bipartite graphs]
{The Ferrers bound for \\spanning trees in bipartite graphs}
\author{Boon Suan Ho}
\address{Department of Mathematics, National University of Singapore}
\email{\href{mailto:hbs@u.nus.edu}{hbs@u.nus.edu}}
\date{}

\begin{document}
\begin{abstract}
We prove Ehrenborg's conjecture that every connected bipartite graph $G$ with 
parts of size $m$ and $n$ has at most $\frac{1}{mn}\prod_{v\in V(G)}\deg(v)$ 
spanning trees, and that equality holds if and only if $G$ is a Ferrers graph. 
The proof is fully formalized in Lean 4.
\end{abstract}

\maketitle

\section{Introduction}
Given a bipartite graph $G=(X\sqcup Y,E)$, its \emph{Ferrers invariant} is defined by
\[
F(G)\coloneqq \frac{1}{|X||Y|}\prod_{v\in V(G)} \deg(v).
\]
Ehrenborg conjectured in 2006 that the number $\tau(G)$ of spanning trees in 
$G$ is at most $F(G)$ for every bipartite $G$ (see \cite{Slone}). In this paper, 
we prove this conjecture and characterize the equality case.
To avoid trivial issues, we will assume that all graphs are connected in what follows.

The origin of the problem lies in the class of \emph{Ferrers graphs}, studied by 
Ehrenborg and van Willigenburg~\cite{EhrenborgVW}. In one convenient formulation, 
a connected bipartite graph is Ferrers if the neighborhoods of the vertices on one 
side can be linearly ordered by inclusion. 

\begin{figure}[htbp]
    \centering
    \includegraphics[width=0.8\textwidth]{FerrersGraph.pdf}
    \caption{The edges of a Ferrers graph are in direct correspondence with 
    the boxes in a Ferrers diagram. (Reproduced from David Eppstein, 
    ``Four open from IPAM''~\cite{EppsteinIPAM}, 2009, CC BY 4.0 license.)}
    \label{fig:ferrers-graph}
\end{figure}

Ehrenborg and van Willigenburg proved that Ferrers graphs satisfy the exact 
product formula
\[
\tau(G)=F(G),
\]
and in fact established a weighted refinement, using ideas from the theory of 
electrical networks \cite{EhrenborgVW}. The Ferrers equality case has since 
been revisited from several directions: Burns gave bijective proofs 
\cite{Burns}, Bapat gave another proof via resistance distance 
\cite{Bapat}, Volkova obtained a proof by inductively computing the 
determinant of a matrix \cite{Volkova}, and Go, Khwa, Luo, and Stamps 
later rederived the Ferrers formula using triangular rank-one 
perturbations of Laplacian matrices \cite{GoKhwaLuoStamps}.

The general conjecture has nevertheless resisted earlier attacks. 
Garrett and Klee~\cite{GarrettKlee} reformulated the conjecture in terms 
of the nonnegativity of a certain homogeneous polynomial, and used this 
reformulation to verify the conjecture when one part has size at most $5$. 
They also proved separately that adding a pendant vertex preserves the 
conjecture, which yields the tree case.
See also Slone's survey~\cite{Slone} and Bapat's exposition \cite{Bapat}. 
Koo's senior thesis \cite{Koo} established further special cases, including 
even cycles and sufficiently dense bipartite graphs with a cut vertex of 
degree~$2$, and also showed that the operation of connecting two graphs by 
a new edge preserves Ferrers-goodness (that is, preserves the bound 
$\tau(G)\le F(G)$). Slone~\cite{Slone} also proved that gluing Ferrers-good graphs 
along a vertex preserves Ferrers-goodness, and consequently that any minimal 
counterexample must be $2$-connected; the same paper recorded computational 
verification up to $13$ vertices and proposed a stronger majorization-based 
statement that was shown there to be false. 
Later work of Slone \cite{Slone2} extended the computational verification 
from $13$ to $17$ vertices.
In 2020, Volkova proved the conjecture for one-side regular bipartite 
graphs~\cite{Volkova}.

At the same time, several complementary reformulations and weaker bounds were 
developed. Bozkurt proved the general bipartite estimate
\[
\tau(G)\le \frac{1}{|E(G)|}\prod_{v\in V(G)} \deg(v),
\]
with equality exactly for complete bipartite graphs \cite{Bozkurt}; this is 
weaker than the Ferrers bound by a factor of $1$ over the bipartite density 
$|E(G)|/(|X||Y|)$. Spectral and majorization reformulations were recorded by 
Slone and Bapat \cite{Slone,Bapat}, and more recently Cherkashin and Prozorov 
showed that a natural polynomial version of Ehrenborg's conjecture is 
equivalent to the original inequality \cite[Theorem 2]{CherkashinProzorov}. 
Consequently, the following is a corollary of our result:

\begin{corollary}
Let $G=(X\sqcup Y,E)$ be a bipartite graph. Define
\[
P_G\bigl((z_v)_{v\in V(G)}\bigr)\coloneqq
\sum_T \prod_{v\in V(G)} z_v^{\deg_T(v)-1},
\]
where the sum ranges over all spanning trees $T$ of $G$ (so $P_G$ is the zero polynomial if $G$ is disconnected). Then, for every choice of nonnegative reals $z_v$ $(v\in V(G))$, we have
\[
P_G\bigl((z_v)_{v\in V(G)}\bigr)
\Bigl(\sum_{x\in X} z_x\Bigr)
\Bigl(\sum_{y\in Y} z_y\Bigr)
\le
\prod_{v\in V(G)} \Bigl(\sum_{u\in N_G(v)} z_u\Bigr).
\]
\end{corollary}

The present argument is linear-algebraic in the spirit of Klee and Stamps
\cite{KleeStampsWeighted,KleeStamps}. Starting from Kirchhoff's matrix-tree
theorem, we take a Schur complement of the Laplacian matrix to reduce the
problem to one side of the bipartition. This yields the matrix
\[
L_X=A-BC^{-1}B^\top,
\]
and turns the spanning-tree bound into a determinant inequality on the
$X$-side. After a rank-one perturbation, the resulting positive definite matrix
\[
M=L_X+\frac{n}{m}J
\]
admits a decomposition as a sum of orthogonal projections indexed by the 
neighborhoods of the vertices of $Y$. This structural observation makes 
Ky Fan's maximum principle directly applicable and produces a majorization 
of the degree sequence of $X$ by the eigenvalues of $M$. The determinant 
bound then follows from the strict Schur-concavity of $\sum\log x_i$, 
and the equality case is recovered from an explicit nonnegative overlap 
defect in the projection traces.

Now let $G=(X\sqcup Y,E)$ be a bipartite graph with $|X|=m$ and $|Y|=n$. Write
\[
D(G)\coloneqq \prod_{v\in V(G)}\deg(v),
\]
where $\deg(v)$ denotes the degree of $v$, and let $\tau(G)$ be the number 
of spanning trees of $G$.
Following \cite{EhrenborgVW}, we call a bipartite graph $G=(X\sqcup Y,E)$ 
a \emph{Ferrers graph} if, after ordering
\[
X=\{x_1,\dots,x_m\},\qquad Y=\{y_1,\dots,y_n\},
\]
there exist integers
\[
m=t_1\ge t_2\ge \cdots \ge t_n\ge 1
\]
such that
\[
N(y_j)=\{x_1,\dots,x_{t_j}\}\qquad (1\le j\le n).
\]
Equivalently, the neighborhoods on one side of the bipartition are 
linearly ordered by inclusion. In terms of the biadjacency matrix, 
this means that after a suitable ordering, every column consists of 
a block of $1$'s followed by $0$'s, and the column heights are weakly 
decreasing; in other words, the pattern of $1$'s forms a Ferrers diagram. 

The main result of this paper is the following:

\begin{theorem}\label{thm:main}
Let $G=(X\sqcup Y,E)$ be a connected bipartite graph with $|X|=m$ and $|Y|=n$. 
Then
\[
\tau(G)\le \frac{1}{mn}\prod_{v\in V(G)}\deg(v).
\]
Moreover, equality holds if and only if $G$ is a Ferrers graph.
\end{theorem}

\section{Determinant reduction}

\subsection*{Standard ingredients}
We will use the following standard results: Kirchhoff's matrix-tree theorem 
from algebraic graph theory \cite[Section~13.2]{GodsilRoyle}, the determinant 
formula for Schur complements from matrix analysis 
\cite[Section~0.8.5]{HornJohnson}, Ky Fan's maximum principle 
\cite[Exercise II.1.13]{Bhatia}, and the basic majorization principle 
for strictly concave functions \cite[Section~3C]{MarshallOlkinArnold}.

\begin{theorem}[Kirchhoff's matrix-tree theorem]\label{thm:kirchhoff}
Let $H$ be a graph with Laplacian matrix $L(H)$. Deleting the same row 
and column from $L(H)$ produces a submatrix whose determinant is independent 
of the deleted vertex and equals $\tau(H)$.
\end{theorem}

\begin{proposition}[Schur complement determinant formula]
\label{prop:schur-complement}
Let $U$ and $Z$ be square matrices, 
and let $V$ and $W$ be compatible rectangular blocks. 
If $Z$ is invertible, then
\[
\det\begin{pmatrix}
U & V\\
W & Z
\end{pmatrix}
=
\det(Z)\,\det(U-VZ^{-1}W).
\]
\end{proposition}

\begin{restatable}[Ky Fan's maximum principle \cite{KyFan}]{theorem}{kyfan}
\label{thm:ky-fan}
Let $S$ be a real symmetric $m\times m$ matrix with eigenvalues
\[
\theta_1\ge \theta_2\ge \cdots \ge \theta_m.
\]
For each $k\in\{1,\dots,m\}$,
\[
\sum_{i=1}^k \theta_i
=
\max\bigl\{\operatorname{tr}(PS):\ P \text{ is an orthogonal projection of rank }k\bigr\}.
\]
In particular, for every rank-$k$ orthogonal projection $P$,
\[
\operatorname{tr}(PS)\le \sum_{i=1}^k \theta_i.
\]
\end{restatable}
\begin{proof}See Appendix \ref{sec:A}.\end{proof}

\begin{definition}\label{def:majorization}
For $x,y\in\RR^m$, write $x^\downarrow$ and $y^\downarrow$ for their coordinates 
rearranged in weakly decreasing order. We say that $x$ \emph{majorizes} $y$ if
\[
\sum_{i=1}^k x_i^\downarrow\ge \sum_{i=1}^k y_i^\downarrow\qquad (1\le k<m)
\]
and
\[
\sum_{i=1}^m x_i=\sum_{i=1}^m y_i.
\]
\end{definition}

\begin{proposition}\label{prop:strict-schur}
If $x,y\in(0,\infty)^m$ and $x$ majorizes $y$, then for every strictly 
concave function $\phi\colon(0,\infty)\to\RR$,
\[
\sum_{i=1}^m \phi(x_i)\le \sum_{i=1}^m \phi(y_i).
\]
If, in addition, $x$ and $y$ are both written in weakly decreasing order, 
then equality holds if and only if $x_i=y_i$ for all $i$.
\end{proposition}

\begin{proof}
This is the standard strict form of Karamata's inequality, or equivalently the 
statement that $\sum_i \phi(x_i)$ is strictly Schur-concave when $\phi$ is 
strictly concave; see \cite[Section 3C]{MarshallOlkinArnold}.
\end{proof}

\subsection*{\texorpdfstring{The Schur complement on $X$}{The Schur complement on X}}
Order the vertices so that those of $X$ come first. 
The Laplacian of $G$ then has block form
\[
L(G)=
\begin{pmatrix}
A & -B\\
-B^\top & C
\end{pmatrix},
\]
where $B\in\{0,1\}^{m\times n}$ is the biadjacency matrix 
(so that $B_{ij}=1$ if $x_i\adj y_j$ and is $0$ otherwise) and
\[
A=\diag(a_1,\dots,a_m),\qquad C=\diag(b_1,\dots,b_n),
\]
with $a_i=\deg(x_i)$ and $b_j=\deg(y_j)$. 
Because $G$ is connected, every $b_j$ is positive, so $C$ is invertible. 
For each $y_j\in Y$, write
\[
T_j\coloneqq N(y_j)\subseteq X,\]
noting that $b_j=|T_j|$. For every nonempty $T\subseteq X$, 
let $\mathbf1_T\in\RR^m$ be its indicator vector and define
\[
D_T\coloneqq \diag(\mathbf1_T),\qquad 
J_T\coloneqq \mathbf1_T\mathbf1_T^{\top},\qquad 
P_T\coloneqq D_T-\frac{1}{|T|}J_T,
\]
together with the subspace
\[
H_T\coloneqq 
\Bigl\{u\in\RR^m:
\ \operatorname{supp}(u)\subseteq T,\ \sum_{x_i\in T}u_i=0\Bigr\}.
\]
For each $u\in\RR^m$, the vector $P_Tu$ is obtained by restricting $u$ to $T$, subtracting the average value on $T$, and then extending by $0$ outside $T$:
\[
(P_Tu)_i=
\begin{cases}
u_i-\bar u_T& \text{if } x_i\in T,\\ 
0& \text{if } x_i\notin T;
\end{cases}
\]
where $\bar u_T\coloneqq \frac{1}{|T|}\sum_{x_r\in T}u_r$. 
In particular, $P_Tu\in H_T$ for every $u$. Conversely, if $w\in H_T$, 
then $\bar w_T=0$ and $P_Tw=w$. Moreover, for every $w\in H_T$,
\[
\langle u-P_Tu,w\rangle
=
\sum_{x_i\in T}\bar u_T\,w_i
=
\bar u_T\sum_{x_i\in T}w_i
=
0.
\]
Thus $P_T$ is the orthogonal projection onto $H_T$. 
In particular, $P_T$ is symmetric positive semidefinite and 
$P_T\mathbf1=0$. Equivalently, it is the Laplacian of the complete graph 
on $T$, divided by $|T|$ and padded with zeroes outside $T$.

Since the $j$th column of $B$ is $\mathbf1_{T_j}$ and 
$C^{-1}=\diag(|T_1|^{-1},\dots,|T_n|^{-1})$, 
we have
\[
BC^{-1}B^\top=\sum_{j=1}^n\frac{1}{|T_j|}\mathbf1_{T_j}\mathbf1_{T_j}^\top
=\sum_{j=1}^n \frac{1}{|T_j|}J_{T_j},
\qquad
A=\sum_{j=1}^n D_{T_j};
\]
and therefore the Schur complement $L_X\coloneqq A-BC^{-1}B^\top$ satisfies
\begin{equation}\label{eq:L-sumP}
L_X
=\sum_{j=1}^n P_{T_j}.
\end{equation}
This is the first key identity of the proof: each vertex $y_j$ contributes 
the orthogonal projection onto the zero-sum subspace $H_{T_j}$ of its neighborhood.

\begin{remark}
One way to interpret $L_X$ is as follows: 
If we assign weights $u\in\mathbb R^m$ to the $X$-vertices 
and $v\in\mathbb R^n$ to the $Y$-vertices, the quadratic form of $L(G)$ is
\[
\begin{pmatrix}u\\v\end{pmatrix}^\top
L(G)\begin{pmatrix}u\\v\end{pmatrix}
=\sum_{x_iy_j\in E}(u_i-v_j)^2,
\]
which we may view as an energy to minimize for fixed $u\in\mathbb R^m$ 
and varying $v\in\mathbb R^n$. Now a vertex $y_j$ with neighborhood 
$T_j\subseteq X$ contributes
\[
\sum_{x_i\in T_j}(u_i-v_j)^2
\]
to this energy, and this contribution is minimized when $v_j$ is the 
average of its neighbors: 
\[
v_j=\bar u_{T_j}\coloneqq\frac{1}{|T_j|}\sum_{x_i\in T_j}u_i.
\]
The minimized energy from the neighborhood of vertex $y_j$ is then
\[
u^\top P_{T_j}u=\sum_{x_i\in T_j}(u_i-\bar u_{T_j})^2,
\]
and the total minimized energy is
\[
u^\top L_Xu
=\sum_{x_iy_j\in E}(u_i-\bar u_{T_j})^2
=\min_{v\in\mathbb R^n}\sum_{x_iy_j\in E}(u_i-v_j)^2.
\]
\end{remark}

\begin{lemma}\label{lem:kernel}
The matrix $L_X$ is symmetric positive semidefinite, satisfies 
$L_X\mathbf1=0$, and has kernel exactly $\langle \mathbf1\rangle$. 
In particular, its eigenvalues are
\[
0=\mu_1<\mu_2\le \cdots \le \mu_m.
\]
\end{lemma}

\begin{proof}
By the discussion above, each $P_T$ is an orthogonal projection, hence symmetric positive semidefinite, and satisfies $P_T\mathbf1=0$. Therefore \eqref{eq:L-sumP} implies that $L_X$ is symmetric positive semidefinite and that
\[
L_X\mathbf1=0.
\]

Now let $u\in\ker L_X$. Since $L_X$ is a sum of positive semidefinite matrices,
\[
0=u^{\top}L_Xu=\sum_{j=1}^n u^{\top}P_{T_j}u=\sum_{j=1}^n \|P_{T_j}u\|^2,
\]
so $P_{T_j}u=0$ for every $j$. By the displayed formula for $P_Tu$, this means that $u$ is constant on each neighborhood $T_j$.

If $x,x'\in X$, connectedness of $G$ gives an alternating path
\[
x=z_0\adj w_1\adj z_1\adj w_2\adj\cdots\adj w_t\adj z_t=x'.
\]
Since both $z_{s-1}$ and $z_s$ lie in $N(w_s)$, the constancy of 
$u$ on each $T=N(w_s)$ implies
\[
u(z_{s-1})=u(z_s)\qquad (1\le s\le t),
\]
and hence $u(x)=u(x')$. Thus $u$ is constant on all of $X$, 
proving $\ker L_X=\langle \mathbf1\rangle$.
\end{proof}

\begin{proposition}\label{prop:reduction}
Let $M\coloneqq L_X+\frac{n}{m}J$, where $J$ is the $m\times m$ all-ones matrix. Then
\begin{equation}\label{eq:tau-detM}
\tau(G)=\frac{\prod_{j=1}^n b_j}{mn}\det M.
\end{equation}
Consequently, the inequality in Theorem~\ref{thm:main} is equivalent to
\begin{equation}\label{eq:det-bound}
\det M\le \prod_{i=1}^m a_i.
\end{equation}
\end{proposition}

\begin{proof}
Write the Laplacian of $G$ in the block form introduced above, 
and remove the row and column corresponding to $x_1$. 
By Kirchhoff's matrix-tree theorem (Theorem~\ref{thm:kirchhoff}),
\[
\tau(G)
=\det\begin{pmatrix}
A_{\hat 1} & -B_{\hat 1}\\
-B_{\hat 1}^{\top} & C
\end{pmatrix}.
\]
Since $C$ is invertible, the Schur complement determinant formula 
(Proposition~\ref{prop:schur-complement}) applied to the $C$-block gives
\[
\tau(G)
=\det(C)\,\det\bigl(A_{\hat 1}-B_{\hat 1}C^{-1}B_{\hat 1}^{\top}\bigr).
\]
By definition, $L_X=A-BC^{-1}B^\top$, so
\[
A_{\hat 1}-B_{\hat 1}C^{-1}B_{\hat 1}^{\top}
=\bigl(A-BC^{-1}B^\top\bigr)_{\hat 1}
=\bigl(L_X\bigr)_{\hat 1}.
\]
Hence
\begin{equation}\label{eq:kirchhoff-schur-step}
\tau(G)=\det(C)\,\det\bigl((L_X)_{\hat 1}\bigr).
\end{equation}

Now let $0,\mu_2,\dots,\mu_m$ be the eigenvalues of $L_X$, as in 
Lemma~\ref{lem:kernel}. Since $L_X$ is symmetric and 
$\ker L_X=\langle \mathbf1\rangle$, we may choose an orthonormal eigenbasis 
whose first vector is $m^{-1/2}\mathbf1$. Let $Q$ be the orthogonal matrix 
whose columns are this basis. Then
\[
L_X=Q\diag(0,\mu_2,\dots,\mu_m)Q^\top.
\]
Since the adjugate matrix (i.e., the transpose of the cofactor matrix) 
satisfies $\operatorname{adj}(SDS^{-1})=S\,\operatorname{adj}(D)\,S^{-1}$, we obtain
\[
\operatorname{adj}(L_X)
=Q\diag\Bigl(\prod_{i=2}^m\mu_i,0,\dots,0\Bigr)Q^\top.
\]
If $u=m^{-1/2}\mathbf1$, then $u=Qe_1$, so
\[
Q\diag\Bigl(\prod_{i=2}^m\mu_i,0,\dots,0\Bigr)Q^\top
=\Bigl(\prod_{i=2}^m\mu_i\Bigr)uu^\top.
\]
Since $uu^\top=\frac1m J$, this gives
\[
\operatorname{adj}(L_X)=\frac1m\Bigl(\prod_{i=2}^m\mu_i\Bigr)J.
\]
So every cofactor of $L_X$ is the same, and in particular
\[
\det\bigl((L_X)_{\hat 1}\bigr)=\frac{1}{m}\prod_{i=2}^m \mu_i.
\]

Finally, $J\mathbf1=m\mathbf1$, while $J$ vanishes on $\mathbf1^\perp$. 
Thus on the line $\langle \mathbf1\rangle$, the matrix
\[
M=L_X+\frac{n}{m}J
\]
acts by $n$, and on $\mathbf1^\perp$ it agrees with $L_X$. 
Hence the eigenvalues of $M$ are
\[
n,\mu_2,\dots,\mu_m,
\]
which are all strictly positive by Lemma~\ref{lem:kernel};
in particular, $M$ is positive definite. So
\[
\det M=n\prod_{i=2}^m\mu_i
=mn\,\det\bigl((L_X)_{\hat 1}\bigr).
\]
Combining this with \eqref{eq:kirchhoff-schur-step} 
and $\det(C)=\prod_{j=1}^n b_j$, we get
\[
\tau(G)
=\det(C)\,\det\bigl((L_X)_{\hat 1}\bigr)
=\Bigl(\prod_{j=1}^n b_j\Bigr)\det\bigl((L_X)_{\hat 1}\bigr)
=\frac{\prod_{j=1}^n b_j}{mn}\det M,
\]
which is \eqref{eq:tau-detM}. 
The equivalence with \eqref{eq:det-bound} is immediate.
\end{proof}


\section{Projection identities}

For every nonempty $T\subseteq X$, define
\[
Q_T\coloneqq P_T+\frac{1}{m}J.
\]
Recall that $H_T\subseteq\RR^m$ was defined in the previous section.

\begin{lemma}\label{lem:Q-proj}
For every nonempty $T\subseteq X$, the matrix $Q_T$ is the 
orthogonal projection onto
\[
H_T\oplus \langle \mathbf1\rangle.
\]
In particular, $Q_T$ has rank $|T|$.
\end{lemma}

\begin{proof}
By the discussion after the definition of $P_T$, the matrix $P_T$ is the orthogonal projection onto $H_T$. Also, $\frac{1}{m}J$ is the orthogonal projection onto $\langle \mathbf1\rangle$. Since $H_T\perp\langle \mathbf1\rangle$, their sum
\[
Q_T=P_T+\frac{1}{m}J
\]
is the orthogonal projection onto $H_T\oplus\langle\mathbf1\rangle$. Its rank is
\[
\dim(H_T)+1=(|T|-1)+1=|T|.
\]
\end{proof}

\begin{lemma}\label{lem:M-sumQ}
One has
\begin{equation}\label{eq:M-sumQ}
M=\sum_{j=1}^n Q_{T_j}.
\end{equation}
\end{lemma}

\begin{proof}
By \eqref{eq:L-sumP},
\[
L_X=\sum_{j=1}^n P_{T_j}.
\]
Adding
\[
\frac{n}{m}J=\sum_{j=1}^n \frac{1}{m}J
\]
termwise gives \eqref{eq:M-sumQ}.
\end{proof}

\begin{remark}
Equation~\eqref{eq:M-sumQ} is the second key identity of the proof.
It rewrites $M$ as a sum of neighborhood projections. Consequently, for
every nonempty subset $I\subseteq X$,
\[
\operatorname{tr}(Q_I M)=\sum_{j=1}^n \operatorname{tr}(Q_IQ_{T_j}).
\]
Lemma~\ref{lem:overlap} will show that each summand is at least
$|I\cap T_j|$, and hence
\[
\operatorname{tr}(Q_I M)\ge \sum_{j=1}^n |I\cap T_j|
= \sum_{x_i\in I} a_i.
\]
This is the bridge from the projection decomposition of $M$ to the
degree sums on the $X$-side used in the majorization argument.
\end{remark}

\begin{lemma}\label{lem:overlap}
For nonempty subsets $I,T\subseteq X$,
\begin{equation}\label{eq:overlap}
\operatorname{tr}(Q_IQ_T)
=
|I\cap T|+\frac{|I\setminus T||T\setminus I|}{|I||T|}
\ge |I\cap T|.
\end{equation}
Equality holds if and only if $I\subseteq T$ or $T\subseteq I$.
\end{lemma}

\begin{proof}
Write
\[
p\coloneqq |I|,\qquad q\coloneqq |T|,\qquad r\coloneqq |I\cap T|.
\]
Because $P_IJ=JP_T=0$ and $J^2=mJ$,
\[
\operatorname{tr}(Q_IQ_T)=\operatorname{tr}(P_IP_T)+1.
\]
Now
\[
P_IP_T=\Bigl(D_I-\frac{1}{p}J_I\Bigr)\Bigl(D_T-\frac{1}{q}J_T\Bigr),
\]
and a direct trace computation gives
\[
\operatorname{tr}(D_ID_T)=r,
\qquad
\operatorname{tr}(D_IJ_T)=r,
\qquad
\operatorname{tr}(J_ID_T)=r,
\qquad
\operatorname{tr}(J_IJ_T)=r^2.
\]
Hence
\[
\operatorname{tr}(P_IP_T)=r-\frac{r}{p}-\frac{r}{q}+\frac{r^2}{pq},
\]
so
\[
\operatorname{tr}(Q_IQ_T)
=
1+r-\frac{r}{p}-\frac{r}{q}+\frac{r^2}{pq}
=
|I\cap T|+\frac{|I\setminus T||T\setminus I|}{|I||T|}.
\]
The error term is nonnegative and vanishes exactly when $r=p$ or $r=q$; 
that is, exactly when $I\subseteq T$ or $T\subseteq I$.
\end{proof}

\begin{remark}
The extra term
\[
\epsov(I,T)\coloneqq \frac{|I\setminus T||T\setminus I|}{|I||T|}
\]
may be viewed as an \emph{overlap defect}: it is symmetric, nonnegative, 
and vanishes exactly when $I$ and $T$ are comparable by inclusion. 
This defect will be the mechanism behind proving the equality case, 
as Ferrers graphs are precisely the connected bipartite graphs such that 
neighborhoods of the vertices on one side can be linearly ordered by inclusion.
\end{remark}

The next proposition isolates the main spectral mechanism of the argument:
the projection decomposition of $M$, together with Ky Fan's maximum principle,
forces the eigenvalues of $M$ to majorize the degree sequence of $X$.

\begin{proposition}\label{prop:majorization}
Let
\[
\lambda_1\ge \lambda_2\ge \cdots \ge \lambda_m
\]
be the eigenvalues of $M$, and relabel the vertices of $X$ so that
\[
a_1\ge a_2\ge \cdots \ge a_m.
\]
For each $k\in\{1,\dots,m\}$, write
\[
[k]\coloneqq \{x_1,\dots,x_k\}.
\]
Then
\begin{equation}\label{eq:partial-sum-ineq}
\sum_{i=1}^k \lambda_i
\ge
\sum_{i=1}^k a_i+\sum_{j=1}^n \epsov([k],T_j)
\ge
\sum_{i=1}^k a_i
\qquad (1\le k<m),
\end{equation}
and
\[
\sum_{i=1}^m \lambda_i=\sum_{i=1}^m a_i.
\]
Consequently, $(\lambda_1,\dots,\lambda_m)$ majorizes $(a_1,\dots,a_m)$.
\end{proposition}

\begin{proof}
First suppose that $1\le k<m$. By Lemma~\ref{lem:Q-proj}, the matrix 
$Q_{[k]}$ is an orthogonal projection of rank $k$. Applying Ky Fan's 
maximum principle (Theorem~\ref{thm:ky-fan}) to the symmetric matrix $M$ gives
\[
\sum_{i=1}^k \lambda_i\ge \operatorname{tr}(Q_{[k]}M).
\]
Using Lemma~\ref{lem:M-sumQ} and Lemma~\ref{lem:overlap},
\[
\operatorname{tr}(Q_{[k]}M)
=\sum_{j=1}^n \operatorname{tr}(Q_{[k]}Q_{T_j})
=\sum_{j=1}^n \bigl|[k]\cap T_j\bigr|
 +\sum_{j=1}^n \epsov([k],T_j).
\]
The first sum is the total degree of the first $k$ vertices of $X$:
\[
\sum_{j=1}^n \bigl|[k]\cap T_j\bigr|=\sum_{i=1}^k a_i,
\]
because every edge incident to one of $x_1,\dots,x_k$ is counted exactly once. 
This proves~\eqref{eq:partial-sum-ineq}.
Now suppose $k=m$. Then we have
\[
\sum_{i=1}^m \lambda_i
=\operatorname{tr}(M)
=\sum_{j=1}^n \operatorname{tr}(Q_{T_j})
=\sum_{j=1}^n |T_j|
=\sum_{i=1}^m a_i,
\]
so it follows that $(\lambda_1,\dots,\lambda_m)$ majorizes $(a_1,\dots,a_m)$.
\end{proof}

\section{Proof of the theorem}

\begin{proof}[Proof of Theorem~\ref{thm:main}]
Let
\[
\lambda_1\ge \lambda_2\ge \cdots \ge \lambda_m
\]
be the eigenvalues of $M$, and relabel the vertices of $X$ so that
\[
a_1\ge a_2\ge \cdots \ge a_m.
\]
For each $k\in\{1,\dots,m\}$, write
\[
[k]\coloneqq \{x_1,\dots,x_k\}.
\]
By Proposition~\ref{prop:majorization}, 
the vector $(\lambda_1,\dots,\lambda_m)$ majorizes 
the degree vector $(a_1,\dots,a_m)$.

By Lemma~\ref{lem:kernel}, $L_X$ has eigenvalues $0,\mu_2,\dots,\mu_m$, 
and since $J$ acts by $m$ on $\langle\mathbf1\rangle$ and by $0$ on 
$\mathbf1^\perp$, $M$ has eigenvalues $n,\mu_2,\dots,\mu_m$, and is 
thus positive definite, so we have $\lambda_i>0$ for every $i$. 
Because $G$ is connected, every vertex degree is positive, so $a_i>0$ 
for every $i$ as well. Thus both vectors lie in $(0,\infty)^m$, and 
Proposition~\ref{prop:strict-schur} applied with the strictly concave 
function $\phi(t)=\log t$ gives
\[
\sum_{i=1}^m \log \lambda_i\le \sum_{i=1}^m \log a_i.
\]
Equivalently,
\[
\det M=\prod_{i=1}^m \lambda_i\le \prod_{i=1}^m a_i.
\]
By Proposition~\ref{prop:reduction},
\[
\tau(G)=\frac{\prod_{j=1}^n b_j}{mn}\det M
\le
\frac{1}{mn}\Bigl(\prod_{i=1}^m a_i\Bigr)\Bigl(\prod_{j=1}^n b_j\Bigr)
=
\frac{D(G)}{mn},
\]
which proves the inequality.

Now suppose equality holds in Theorem~\ref{thm:main}. 
Then equality holds in \eqref{eq:det-bound}, so
\[
\sum_{i=1}^m \log \lambda_i=\sum_{i=1}^m \log a_i.
\]
We already know from Proposition~\ref{prop:majorization} that 
$(\lambda_1,\dots,\lambda_m)$ majorizes $(a_1,\dots,a_m)$, and both 
vectors were arranged in weakly decreasing order at the start of the proof. 
Therefore the equality statement in Proposition~\ref{prop:strict-schur} 
yields coordinatewise equality,
\[
\lambda_i=a_i\qquad (1\le i\le m).
\]
Returning to \eqref{eq:partial-sum-ineq}, we get
\[
\sum_{j=1}^n \epsov([k],T_j)=0
\qquad (1\le k<m).
\]
Each summand is nonnegative, so every summand vanishes:
\[
\epsov([k],T_j)=0
\qquad\text{for all }j\text{ and }1\le k<m.
\]
By Lemma~\ref{lem:overlap}, each $T_j$ is comparable by inclusion with every 
initial segment $[k]$.

We now show that this forces $T_j$ itself to be an initial segment. Fix $j$ and let
\[
t\coloneqq \max\{i:\ x_i\in T_j\}.
\]
Then $T_j\subseteq [t]$. If $t=1$, then $T_j\neq\varnothing$ implies $T_j=[1]$. 
Assume from now on that $t\ge 2$. Since $x_t\in T_j$, the inclusion 
$T_j\subseteq [t-1]$ is impossible. Comparability of $T_j$ with $[t-1]$ 
therefore gives
\[
[t-1]\subseteq T_j.
\]
Together with $x_t\in T_j$, this yields $T_j=[t]$. Thus every neighborhood of 
a $Y$-vertex is an initial segment of $X$. Reordering the vertices of $Y$ so 
that these initial segments appear in decreasing size, 
we conclude that $G$ is Ferrers.

Conversely, suppose $G$ is Ferrers. Then after ordering $X$ and $Y$ we have
\[
T_j=[t_j]\qquad\text{with}\qquad t_1\ge\cdots\ge t_n\ge 1.
\]
The subspaces
\[
\operatorname{im}(Q_{[1]})\subset \operatorname{im}(Q_{[2]})
\subset \cdots \subset \operatorname{im}(Q_{[m]})=\RR^m
\]
form a complete flag, and $\dim \operatorname{im}(Q_{[t]})=t$ by 
Lemma~\ref{lem:Q-proj}. Choose an orthonormal basis $u_1,\dots,u_m$ adapted 
to this flag, meaning that $u_1,\dots,u_t$ span $\operatorname{im}(Q_{[t]})$ 
for every~$t$. In this basis each $Q_{[t]}$ is diagonal with exactly the 
first~$t$ diagonal entries equal to~$1$. Hence
\[
M=\sum_{j=1}^n Q_{T_j}=\sum_{j=1}^n Q_{[t_j]}
\]
is diagonal, with diagonal entries
\[
\#\{j:\ t_j\ge 1\},\ \#\{j:\ t_j\ge 2\},\ \dots,\ \#\{j:\ t_j\ge m\}.
\]
But $\#\{j: t_j\ge r\}$ is exactly the degree of $x_r$, since $x_r$ is 
adjacent precisely to those~$y_j$ with $t_j\ge r$. Therefore the diagonal 
entries are $a_1,\dots,a_m$, so
\[
\det M=\prod_{i=1}^m a_i.
\]
Proposition~\ref{prop:reduction} now gives equality in Theorem~\ref{thm:main}.
\end{proof}

\section*{Acknowledgements}
The author used GPT-5.4 Pro during the development of this work to 
explore possible proof strategies, test intermediate formulations, 
and assist with exposition. All mathematical arguments and claims in 
the final manuscript were independently verified by the author, who 
takes full responsibility for the paper. 
The Lean~4 formalization is available from
\url{https://github.com/boonsuan/FerrersBound}.

\begin{thebibliography}{99}

\bibitem{Bapat}
Ravindra B.~Bapat,
Maximizing spectral radius and number of spanning trees in bipartite graphs,
in \emph{Mathematical Programming and Game Theory},
Indian Statistical Institute Series, Springer, Singapore, 2018, pp.~33--48.
\href{https://doi.org/10.1007/978-981-13-3059-9_2}{\nolinkurl{doi:10.1007/978-981-13-3059-9\_2}}.

\bibitem{Bhatia}
Rajendra Bhatia,
\emph{Matrix Analysis},
Graduate Texts in Mathematics, vol.~169, Springer-Verlag, New York, 1997.
\href{https://doi.org/10.1007/978-1-4612-0653-8}{\nolinkurl{doi:10.1007/978-1-4612-0653-8}}.

\bibitem{Bozkurt}
\c{S}erife B.~Bozkurt,
Upper bounds for the number of spanning trees of graphs,
\emph{Journal of Inequalities and Applications} \textbf{2012} (2012), article 269.
\href{https://doi.org/10.1186/1029-242X-2012-269}{\nolinkurl{doi:10.1186/1029-242X-2012-269}}.

\bibitem{Burns}
Jason Burns,
Bijective proofs for ``Enumerative properties of Ferrers graphs'',
\href{https://arxiv.org/abs/math/0312282}{\nolinkurl{arXiv:math/0312282}}
\texttt{[math.CO]}, 2003.

\bibitem{CherkashinProzorov}
Danila Cherkashin and Pavel Prozorov,
Around the number of trees in distance-hereditary graphs,
\emph{Discrete Mathematics} \textbf{349} (2026), no.~6, 115010.
\href{https://doi.org/10.1016/j.disc.2026.115010}{\nolinkurl{doi:10.1016/j.disc.2026.115010}}.

\bibitem{EhrenborgVW}
Richard Ehrenborg and Stephanie van Willigenburg,
Enumerative properties of Ferrers graphs,
\emph{Discrete \& Computational Geometry} \textbf{32} (2004), no.~4, 481--492.
\href{https://doi.org/10.1007/s00454-004-1135-1}{\nolinkurl{doi:10.1007/s00454-004-1135-1}}.

\bibitem{EppsteinIPAM}
David Eppstein,
\emph{Four open from IPAM},
blog post, Oct.~24, 2009.
Available at \url{https://11011110.github.io/blog/2009/10/24/four-open-from.html}.

\bibitem{KyFan}
Ky Fan,
On a theorem of Weyl concerning eigenvalues of linear transformations. I,
\emph{Proceedings of the National Academy of Sciences U.S.A.} \textbf{35} (1949), 652--655.
\href{https://doi.org/10.1073/pnas.35.11.652}{\nolinkurl{doi:10.1073/pnas.35.11.652}}.

\bibitem{GarrettKlee}
Fintan Garrett and Steven Klee,
Upper bounds for the number of spanning trees in a bipartite graph,
preprint, 2014.
Available at \url{http://fac-staff.seattleu.edu/klees/web/bipartite.pdf}.

\bibitem{GodsilRoyle}
Chris Godsil and Gordon Royle,
\emph{Algebraic Graph Theory},
Graduate Texts in Mathematics, vol.~207, Springer-Verlag, New York, 2001.
\href{https://doi.org/10.1007/978-1-4613-0163-9}{\nolinkurl{doi:10.1007/978-1-4613-0163-9}}.

\bibitem{GoKhwaLuoStamps}
Christian Go, Zhong Xuan Khwa, Xinyu Luo, and Matthew T. Stamps,
Spanning tree enumeration and nearly triangular graph Laplacians,
\emph{Discrete Applied Mathematics} \textbf{320} (2022), 356--369.
\href{https://doi.org/10.1016/j.dam.2022.04.022}{\nolinkurl{doi:10.1016/j.dam.2022.04.022}}.

\bibitem{HornJohnson}
Roger A.~Horn and Charles R.~Johnson,
\emph{Matrix Analysis},
2nd ed., Cambridge University Press, Cambridge, 2012.
\href{https://doi.org/10.1017/CBO9781139020411}{\nolinkurl{doi:10.1017/CBO9781139020411}}.

\bibitem{KleeStampsWeighted}
Steven Klee and Matthew T.~Stamps,
Linear algebraic techniques for weighted spanning tree enumeration,
\emph{Linear Algebra and its Applications} \textbf{582} (2019), 391--402.
\href{https://doi.org/10.1016/j.laa.2019.08.009}{\nolinkurl{doi:10.1016/j.laa.2019.08.009}}.

\bibitem{KleeStamps}
Steven Klee and Matthew T.~Stamps,
Linear algebraic techniques for spanning tree enumeration,
\emph{American Mathematical Monthly} \textbf{127} (2020), no.~4, 297--307.
\href{https://doi.org/10.1080/00029890.2020.1708171}{\nolinkurl{doi:10.1080/00029890.2020.1708171}}.

\bibitem{Koo}
Cheng Wai Koo,
\emph{A Bound on the Number of Spanning Trees in Bipartite Graphs},
senior thesis, Harvey Mudd College, 2016.
Available at \url{https://scholarship.claremont.edu/hmc_theses/73/}.

\bibitem{MarshallOlkinArnold}
Albert W.~Marshall, Ingram Olkin, and Barry C.~Arnold,
\emph{Inequalities: Theory of Majorization and Its Applications},
2nd ed., Springer, New York, 2011.
\href{https://doi.org/10.1007/978-0-387-68276-1}{\nolinkurl{doi:10.1007/978-0-387-68276-1}}.

\bibitem{Slone}
M.~L.~E.~Slone,
A conjectured bound on the spanning tree number of bipartite graphs,
\href{https://arxiv.org/abs/1608.01929}{\nolinkurl{arXiv:1608.01929}}
\texttt{[math.CO]}, 2016.

\bibitem{Slone2}
M.~L.~E.~Slone,
\emph{Computer verification that no small counterexamples to the Ferrers bound conjecture exist},
GitHub repository, 2025.
Available at \url{https://github.com/cokernel/ferrers-bound-checking}.

\bibitem{Volkova}
Albina Volkova, 
On the number of spanning trees in bipartite graphs,
\href{https://arxiv.org/abs/2009.06688}{\nolinkurl{arXiv:2009.06688}}
\texttt{[math.CO]}, 2020.

\end{thebibliography}

\newpage\appendix 
\section{Proof of Ky Fan's maximum principle} \label{sec:A}
For convenience, we provide a relatively self-contained proof of Ky Fan's 
maximum principle, in the form that we use:
\kyfan*   
\begin{proof}
Let $u_1,\dots,u_m$ be an orthonormal basis of eigenvectors of $S$,
with
\[
Su_i=\theta_i u_i\qquad (1\le i\le m).
\]
Fix a rank-$k$ orthogonal projection $P$, and define
\[
d_i\coloneqq\langle Pu_i,u_i\rangle \qquad (1\le i\le m).
\]
Since $P$ is an orthogonal projection, $P=P^\top=P^2$, and therefore
\[
d_i=\langle P^2u_i,u_i\rangle=\langle Pu_i,Pu_i\rangle=\|Pu_i\|^2.
\]
Hence
\[
0\le d_i\le 1 \qquad (1\le i\le m).
\]
Also, because $(u_i)_{i=1}^m$ is an orthonormal basis,
\[
\sum_{i=1}^m d_i
=\sum_{i=1}^m \langle Pu_i,u_i\rangle
=\operatorname{tr}(P)
=k.
\]

Now compute $\operatorname{tr}(PS)$ in the eigenbasis $(u_i)$:
\[
\operatorname{tr}(PS)
=\sum_{i=1}^m \langle PSu_i,u_i\rangle
=\sum_{i=1}^m \theta_i \langle Pu_i,u_i\rangle
=\sum_{i=1}^m \theta_i d_i.
\]
Therefore
\begin{align*}
\operatorname{tr}(PS)
&=\sum_{i=1}^k \theta_i d_i+\sum_{i=k+1}^m \theta_i d_i
\le\sum_{i=1}^k \theta_i d_i+\theta_k\sum_{i=k+1}^m d_i\\
&\qquad=\sum_{i=1}^k \theta_i d_i+\theta_k\sum_{i=1}^k(1-d_i)
=\sum_{i=1}^k \bigl(\theta_i d_i+\theta_k(1-d_i)\bigr)\\
&\qquad\le\sum_{i=1}^k \bigl(\theta_i d_i+\theta_i(1-d_i)\bigr)
=\sum_{i=1}^k \theta_i,
\end{align*}
where in the first line we used $\theta_i\le \theta_k$ for $i>k$, in the
second we used $\sum_{i=1}^m d_i=k$, and in the third we used
$\theta_k\le \theta_i$ for $i\le k$.

Thus every rank-$k$ orthogonal projection $P$ satisfies
\[
\operatorname{tr}(PS)\le \sum_{i=1}^k \theta_i.
\]

To see that equality is attained, let $P_*$ be the orthogonal
projection onto $\operatorname{span}\{u_1,\dots,u_k\}$. Then
$P_*u_i=u_i$ for $i\le k$ and $P_*u_i=0$ for $i>k$, so
\[
\operatorname{tr}(P_*S)
=\sum_{i=1}^m \theta_i \langle P_*u_i,u_i\rangle
=\sum_{i=1}^k \theta_i.
\]
Hence
\[
\max\bigl\{\operatorname{tr}(PS): P \text{ is an orthogonal projection of rank }k\bigr\}
=\sum_{i=1}^k \theta_i,
\]
which proves the theorem.
\end{proof}
\end{document}
